\section{What the data show (not a model claim)}
\label{sec:data}

Same hike windows as Gate~0: three-month means around base/peak, Pink
Sheet real 2010 dollars~\citep{worldbank_pink}.

\begin{center}
\begin{tabular}{l l rrr}
\toprule
Crisis & Window & Wheat & Rice & Maize \\
\midrule
2007/08 & 2006-06 \(\pm 1\) \(\to\) 2008-03 \(\pm 1\)
  & \(\times 1.82\) & \(\times 1.84\) & \(\times 1.84\) \\
2010/11 & 2009-06 \(\pm 1\) \(\to\) 2011-02 \(\pm 1\)
  & \(\times 1.16\) & \(\times 0.79\) & \(\times 1.44\) \\
\bottomrule
\end{tabular}
\end{center}

Calendar peaks versus the same June-2006 baseline are higher than the
fixed March-2008 window captures: wheat \(\times 1.95\) (March 2008), rice
\(\times 2.54\) (April 2008), maize \(\times 2.26\) (June 2008). Wheat was
already \(\times 1.5\)--\(1.8\) from September 2007 through January 2008
while rice stayed near baseline. The largest rice month-to-month jump was
April 2008, while wheat was already falling.

AMIS rice measures (Vietnam prohibition 2007-07, India from 2007-10,
second wave 2008-02/04) line up with the rice cliff~\citep{amis_oecd}.
\textbf{2007/08 rice is primarily an own-ban episode.} Matching that spike
is not a clean substitution validation target---the model already feeds
rice AMIS. \textbf{2010/11 is the contrast:} wheat and maize up, rice
down. A pass must not invent a rice hike.

Fair empirical bar: wheat hike and timing in both crises; maize
co-movement as the least-confounded spillover; rice 2008 = own-policy
(hard bar 4); rice 2010 = no fake co-spike. Do not require all three
series to look equally pretty.
